\clearpage
\begin{trapbox}
  \textbf{\textcolor{CTrap}{易错一：把“奇函数”误记成“一定经过原点”}}
\begin{center}\begin{tikzpicture}[font=\small,>=stealth,line join=round]
\draw[gray!40,rounded corners=3pt] (0,0) rectangle (14.2,8.5);\node at (7.1,7.9) {易错一：关于原点对称，不等于图像一定包含原点};
\node at (3.5,6.85){$y=x^3$：$0$ 可以代入};\node at (10.6,6.85){$y=1/x$：$0$ 不能代入};\begin{scope}[shift={(3.5,3.7)},x=1cm,y=1cm]\draw[->,gray] (-2.2,0)--(2.2,0) node[right]{$x$};\draw[->,gray] (0,-2.2)--(0,2.2) node[above]{$y$};\draw[blue,thick,domain=-1.25:1.25,samples=150] plot(\x,{(\x)^3});\fill[orange!85!black] (0,0) circle(2.2pt) node[below right=5pt] {$O$};\end{scope}\begin{scope}[shift={(10.5,3.7)},x=1cm,y=1cm]\draw[->,gray] (-2.2,0)--(2.2,0) node[right]{$x$};\draw[->,gray] (0,-2.2)--(0,2.2) node[above]{$y$};\draw[blue,thick,domain=-2:-0.5,samples=150] plot(\x,{1/\x});\draw[blue,thick,domain=0.5:2,samples=150] plot(\x,{1/\x});\fill[orange!85!black] (1,1) circle(2.2pt) node[above right=5pt] {$(1,1)$};\fill[orange!85!black] (-1,-1) circle(2.2pt) node[below left=5pt] {$(-1,-1)$};\node[below right=5pt] at (0,0){$O$};\end{scope}\node at (3.5,1.25){是奇函数，且 $f(0)=0$};\node at (10.6,1.25){是奇函数，但 $f(0)$ 不存在};
\node[text width=12.8cm,align=center] at (7.1,.65) {两图横纵同尺度；右图原点仅是坐标轴交点，不是函数图像上的点。};
\end{tikzpicture}\end{center}



  \textbf{错误说法：}奇函数一定经过原点，所以一定有 $f(0)=0$。

  \textbf{问题在哪里？}奇函数的定义域可以不包含 $0$。
  如 $f(x)=\dfrac1x$ 是奇函数，但 $f(0)$ 无意义，图像上也没有原点这个点。

  \textbf{正确说法：}\textbf{奇函数只要在 $x=0$ 处有定义，就必经过原点。}

  \medskip
  
\end{trapbox}
\clearpage
\begin{trapbox}
\textbf{\textcolor{CTrap}{易错二：把“$f(0)=0$”当成“奇函数”的证明}}
\begin{center}\begin{tikzpicture}[font=\small,>=stealth,line join=round]
\draw[gray!40,rounded corners=3pt] (0,0) rectangle (14.2,8.5);\node at (7.1,7.9) {易错二：同样经过原点，只有一幅是奇函数图像};
\node at (3.5,6.85){$y=x^3$：原点对称};\node at (10.6,6.85){$y=x^2$：仅过原点还不够};\begin{scope}[shift={(3.5,3.3)},x=1cm,y=1cm]\draw[->,gray] (-2.1,0)--(2.1,0) node[right]{$x$};\draw[->,gray] (0,-1.9)--(0,2.8) node[above]{$y$};\draw[blue,thick,domain=-1.2:1.35,samples=150] plot(\x,{(\x)^3});\fill[orange!85!black] (1,1) circle(2.2pt) node[above right=5pt] {$(1,1)$};\fill[orange!85!black] (-1,-1) circle(2.2pt) node[below left=5pt] {$(-1,-1)$};\fill[orange!85!black] (0,0) circle(2.2pt);\end{scope}\begin{scope}[shift={(10.5,3.3)},x=1cm,y=1cm]\draw[->,gray] (-2.1,0)--(2.1,0) node[right]{$x$};\draw[->,gray] (0,-1.9)--(0,2.8) node[above]{$y$};\draw[blue,thick,domain=-1.6:1.6,samples=150] plot(\x,{(\x)^2});\fill[orange!85!black] (1,1) circle(2.2pt) node[above right=5pt] {$(1,1)$};\fill[orange!85!black] (-1,1) circle(2.2pt) node[above left=5pt] {$(-1,1)$};\draw[red,fill=white] (-1,-1) circle(2.5pt);\node[red,right=5pt] at (.25,-1.2){缺少 $(-1,-1)$};\fill[orange!85!black] (0,0) circle(2.2pt);\end{scope}
\node[text width=12.8cm,align=center] at (7.1,.65) {空心红点只表示应有但缺少的对称点，不属于 $y=x^2$ 的图像。};
\end{tikzpicture}\end{center}



  \textbf{错误想法：}已知 $f(0)=0$，图像经过原点，所以函数是奇函数。

  \textbf{反例：}$f(x)=x^2$ 满足 $f(0)=0$，但
  \[
    f(1)=1,\qquad f(-1)=1\ne-1=-f(1).
  \]
  因此它不是奇函数。
  \textbf{正确方法：}要证明奇函数，必须检查定义域关于 $0$ 成对出现，
  并验证每对相反数都满足 $f(-x)=-f(x)$。

  \medskip
  
\end{trapbox}
\clearpage
\begin{trapbox}
\textbf{\textcolor{CTrap}{易错三：求出了一个参数，就以为已经检查完奇函数}}
\begin{center}\begin{tikzpicture}[font=\small,>=stealth,line join=round]
\draw[gray!40,rounded corners=3pt] (0,0) rectangle (14.2,8.5);\node at (7.1,7.9) {易错三：消去常数项，还要消去偶次项};
\node at (3.5,6.85){$y=x^3+x^2$：$r=0$，仍不够};\node at (10.6,6.85){$y=x^3$：还需让 $p=0$};\begin{scope}[shift={(3.5,3.3)},x=1cm,y=1cm]\draw[->,gray] (-2,0)--(2,0) node[right]{$x$};\draw[->,gray] (0,-1.6)--(0,3.0) node[above]{$y$};\draw[blue,thick,domain=-1.6:1.1,samples=150] plot(\x,{(\x)^3+(\x)^2});\fill[orange!85!black] (1,2) circle(2.2pt) node[above right=5pt] {$(1,2)$};\fill[orange!85!black] (-1,0) circle(2.2pt) node[below left=5pt] {$(-1,0)$};\fill[orange!85!black] (0,0) circle(2.2pt);\end{scope}\begin{scope}[shift={(10.5,3.3)},x=1cm,y=1cm]\draw[->,gray] (-2,0)--(2,0) node[right]{$x$};\draw[->,gray] (0,-1.6)--(0,3.0) node[above]{$y$};\draw[blue,thick,domain=-1.1:1.4,samples=150] plot(\x,{(\x)^3});\fill[orange!85!black] (1,1) circle(2.2pt) node[above right=5pt] {$(1,1)$};\fill[orange!85!black] (-1,-1) circle(2.2pt) node[below left=5pt] {$(-1,-1)$};\fill[orange!85!black] (0,0) circle(2.2pt);\end{scope}
\node[text width=12.8cm,align=center] at (7.1,.65) {图取 $q=0$，左图 $p=1,r=0$；右图 $p=r=0$。两图同尺度。};
\end{tikzpicture}\end{center}



  例如已知 $f(x)=x^3+ax^2+bx+c$ 为奇函数，
  由 $f(0)=0$ 只得到 $c=0$。
  \textbf{还要}使用 $f(-x)=-f(x)$ 得到 $a=0$。
  特别是题目让你“求使函数为奇函数的全部参数”时，
  必须把求出的参数代回去验证，而不能只满足 $f(0)=0$ 就停笔。

  \medskip
  
\end{trapbox}
\clearpage
\begin{trapbox}
\textbf{\textcolor{CTrap}{易错四：忘记检查“相反数有没有一起进入定义域”}}
\begin{center}\begin{tikzpicture}[font=\small,>=stealth,line join=round]
\draw[gray!40,rounded corners=3pt] (0,0) rectangle (14.2,8.5);\node at (7.1,7.9) {易错四：同一个表达式 $y=x$，换了定义域就不同};
\node at (3.5,6.85){$D=[0,1]$：不是奇函数};\node at (10.6,6.85){$D=[-1,1]$：是奇函数};\begin{scope}[shift={(3.5,3.7)},x=1.5cm,y=1.5cm]\draw[->,gray] (-1.55,0)--(1.55,0) node[right]{$x$};\draw[->,gray] (0,-1.55)--(0,1.55) node[above]{$y$};\draw[blue,very thick] (0,0)--(1,1);\draw[red,dashed] (-1,-1)--(0,0);\fill[orange!85!black] (0,0) circle(2.2pt) node[below right=5pt] {$O$};\fill[orange!85!black] (1,1) circle(2.2pt) node[above right=5pt] {$(1,1)$};\draw[red,fill=white] (-1,-1) circle(2.5pt);\node[red,below left=5pt] at (-1,-1){$-1\notin D$};\end{scope}\begin{scope}[shift={(10.5,3.7)},x=1.5cm,y=1.5cm]\draw[->,gray] (-1.55,0)--(1.55,0) node[right]{$x$};\draw[->,gray] (0,-1.55)--(0,1.55) node[above]{$y$};\draw[blue,very thick] (-1,-1)--(1,1);\fill[orange!85!black] (1,1) circle(2.2pt) node[above right=5pt] {$(1,1)$};\fill[orange!85!black] (-1,-1) circle(2.2pt) node[below left=5pt] {$(-1,-1)$};\fill[orange!85!black] (0,0) circle(2.2pt);\end{scope}
\node[text width=12.8cm,align=center] at (7.1,.65) {两图同尺度；左图红虚线是缺失的对称像，不属于给定函数。};
\end{tikzpicture}\end{center}



  函数 $f(x)=x$ 如果只规定在 $[0,1]$ 上取值，
  虽然 $f(0)=0$，但 $x=1$ 能取而 $x=-1$ 不能取，
  因此按高中奇函数的定义，它\textbf{不是奇函数}。

  \textbf{记住：}不是只看函数的表达式，还要看题目给出的定义域。

  \medskip
  
\end{trapbox}
\clearpage
\begin{trapbox}
\textbf{\textcolor{CTrap}{易错五：把奇函数和偶函数在 $0$ 处的性质混在一起}}
\begin{center}\begin{tikzpicture}[font=\small,>=stealth,line join=round]
\draw[gray!40,rounded corners=3pt] (0,0) rectangle (14.2,8.5);\node at (7.1,7.9) {易错五：偶函数在零处的值，可以不是零};
\node at (3.5,6.85){奇函数 $y=x^3$};\node at (10.6,6.85){偶函数 $y=x^2+2$};\begin{scope}[shift={(3.5,3.1)},x=0.95cm,y=0.75cm]\draw[->,gray] (-2.2,0)--(2.2,0) node[right]{$x$};\draw[->,gray] (0,-1.5)--(0,4.1) node[above]{$y$};\draw[blue,thick,domain=-1.1:1.55,samples=150] plot(\x,{(\x)^3});\fill[orange!85!black] (0,0) circle(2.2pt) node[below right=5pt] {$(0,0)$};\end{scope}\begin{scope}[shift={(10.5,3.1)},x=0.95cm,y=0.75cm]\draw[->,gray] (-2.2,0)--(2.2,0) node[right]{$x$};\draw[->,gray] (0,-1.5)--(0,4.1) node[above]{$y$};\draw[blue,thick,domain=-1.4:1.4,samples=150] plot(\x,{(\x)^2+2});\fill[orange!85!black] (0,2) circle(2.2pt) node[below right=5pt] {$(0,2)$};\fill[orange!85!black] (-1,3) circle(2.2pt);\fill[orange!85!black] (1,3) circle(2.2pt);\draw[gray,dashed] (-1,3)--(1,3);\end{scope}
\node[text width=12.8cm,align=center] at (7.1,.65) {两图同尺度；左图要求相反数配对，右图只要求左右高度相同。};
\end{tikzpicture}\end{center}



  奇函数：若 $0$ 有定义，则 $f(0)=-f(0)$，所以 $f(0)=0$。

  偶函数：若 $0$ 有定义，则只有 $f(0)=f(0)$，不能确定数值；
  比如 $x^2+2$ 是偶函数，$f(0)=2$。

  \medskip
  
\end{trapbox}
\clearpage
\begin{trapbox}
\textbf{\textcolor{CTrap}{易错六：把“原点处的零点”误当成“唯一零点”}}
\begin{center}\begin{tikzpicture}[font=\small,>=stealth,line join=round]
\draw[gray!40,rounded corners=3pt] (0,0) rectangle (14.2,8.5);\node at (7.1,7.9) {易错六：原点必为零点，不代表它是唯一零点};
\begin{scope}[shift={(4.2,3.8)},x=1.8cm,y=1.5cm]\draw[->,gray] (-1.9,0)--(1.9,0) node[right]{$x$};\draw[->,gray] (0,-1.7)--(0,1.7) node[above]{$y$};\draw[blue,thick,domain=-1.4:1.4,samples=150] plot(\x,{(\x)^3-\x});\fill[orange!85!black] (-1,0) circle(2.2pt) node[below=5pt] {$-1$};\fill[orange!85!black] (0,0) circle(2.2pt) node[above right=8pt] {$0$};\fill[orange!85!black] (1,0) circle(2.2pt) node[below=5pt] {$1$};\draw[gray,dotted] (-1.2,-1.15)--(1.2,-1.15);\draw[<->,orange!85!black] (-1,-1.05)--(1,-1.05);\end{scope}\node[align=left,anchor=west] at (8.4,4.1){$f(x)=x^3-x$\\[11pt]$=x(x-1)(x+1)$\\[11pt]零点：$-1,0,1$\\[11pt]两侧零点成对};
\node[text width=12.8cm,align=center] at (7.1,.65) {曲线直接穿过三个标出的零点；$f(-x)=-f(x)$。};
\end{tikzpicture}\end{center}



  例如 $f(x)=x^3-x$ 是奇函数，
  \[
    f(x)=x(x-1)(x+1),
  \]
  零点有 $-1,0,1$ 三个。
  \textbf{正确理解：}$0$ 在定义域内时必定是零点，但别的零点仍可以成对出现。

  \medskip
  \textbf{\textcolor{CTrap}{考试前最后自查}}

  \begin{itemize}[leftmargin=*, itemsep=0.4em, topsep=0.3em]
    \item $0$ 能不能代入？不能，就不要写 $f(0)$。
    \item 题目是\textbf{已知}奇函数，还是让你\textbf{判断}奇函数？两者做法不同。
    \item 用 $f(0)=0$ 求出参数后，是否还需检查其他位置满足奇函数定义？
  \end{itemize}
\end{trapbox}
